% POST-CHECKPOINT05 WORK IN PROGRESS.
% Reserved map: B.META.22 -> B.148, B.META.23 -> B.149.
% Not an input of the frozen checkpoint05 paper.  No promotion is claimed.
% Existing B146--147 files and their metadata have not been changed.

\subsection{Discrete mixed-characteristic bases without a residue finiteness hypothesis}

\begin{claim}{B.META.22}{Open}{Draft: deep units and a principal-open evaluation detector over an arbitrary mixed-characteristic DVR}
Let $D/\mathbb Q$ be finitely generated and discretely valued
of mixed characteristic $(0,p)$, let $M$ be its completion,
and let $\Omega$ be the algebraic closure of $D$ in a fixed
algebraic closure of $M$.
There is no finite-generation or perfectness assumption on
the residue field of $D$.
Completion identifies $D^h=M\cap\Omega$ and gives an
equivalence on finite separable extensions.  For every
finite extension $L/M$, the field $L\cap\Omega$ is a finite
extension of $D^h$, is dense in $L$, and contains an
$M$-basis of $L$.
For every such complete field $L$ there is a canonical
injective, finite-extension-compatible and Galois-equivariant
map
\[
 \iota_L:L\lhook\joinrel\longrightarrow H^1(G_L,\mathbb Q_p(1)),
 \qquad \iota_L(x)=p^{-r}\kappa_L(\exp(p^r x))\quad(r\gg0).
                                                               \tag{META.AD1}
\]
The inertia character $\chi_p|_{I_M}$ has infinite image,
restriction from $G_M$ to $I_M$ is injective in rational
$p$-Kummer cohomology, and every successive closed derived
quotient of $I_M$ is nonzero.

Here is a local evaluation consequence which does not
require a description of the whole Kummer kernel.
Let $M_0/M$ be finite, let
$N\triangleleft G_{M_0}$ be closed with $\chi_p(N)=1$,
and set $E_0=\Omega^N$.
Suppose that $L/M_0$ is finite Galois and the image of
$N$ in $\operatorname{Gal}(L/M_0)$ is nontrivial.
Let $E/E_0$ be finite, let $V$ be a nonempty principal
open of $\mathbb A^n_E$, and let
$\eta\in H^1(V,\mathbb Q_p(1))$.
Suppose that its evaluation at every point in a nonempty
$L$-open ball of algebraic points
$(L\cap\Omega)^n\cap V$ is zero after the finite scalar
extension containing the coordinates of that point.
Then $\eta=0$.
The same assertion holds if this evaluation hypothesis
is imposed after one fixed finite extension containing
$E$ and $L\cap\Omega$.
\end{claim}
\begin{proof}
All valuations in estimates below are normalized by $v(p)=1$.
We first give the residue-independent ingredients explicitly.
Temkin, Theorem~1.9.5(ii), \S1.9.7 and Lemma~1.9.8,
identifies the henselization of a real-valued field with
its relative separable closure in its completion
\citep{AN_TemkinValuations}.
Thus $D^h=M\cap\Omega$.
Stacks Tag~09EJ, Lemma~15.115.2, whose proof applies to
any DVR, descends every finite separable extension of
the completion \citep{AN_Stacks09EJ}.
Completion preserves degrees: $D^h$ is relatively
algebraically closed in $M$, so is linearly disjoint
from $M$ in every finite Galois extension.
Every finite extension of $D^h$ is henselian and relatively
algebraically closed in its completion.  Its embeddings
preserve the unique extended valuation and extend to
the completions; conversely an embedding of completions
preserves the relative algebraic parts.
This proves full faithfulness as well as essential
surjectivity.  A basis of a finite extension of $D^h$
remains a basis on completion.  These facts prove all
the assertions about $L\cap\Omega$ and identify the
absolute Galois groups and inertia groups by restriction.

For the logarithmic assertion use integer normalization
$v_L(L^\times)=\mathbb Z$ temporarily and write $e_L=v_L(p)$.
Choose $m>e_L/(p-1)$, and write $U_j=1+\mathfrak m_L^j$.
The logarithm and exponential series give inverse
isomorphisms $U_m\simeq\mathfrak m_L^m$.
Indeed $v_L(j!)\le e_Lj/(p-1)$ proves exponential
convergence on this ball, the first term strictly
dominates the subsequent terms, and the formal inverse
identities converge there.
For $u\in U_1$ and $a=v_L(u-1)\ge1$,
\[
             v_L(u^p-1)\ge\min(pa,a+e_L)>a.
\]
Hence $U_1^{p^m}\subset U_m$.
If $z^{p^n}\in U_m$, valuation gives $v_L(z)=0$ and
reduction gives $\overline z^{p^n}=1$.  In any field
of characteristic $p$, perfect or not, the latter
forces $\overline z=1$.
For $n\ge m$ this proves
\[
 U_m^{p^n}\subset U_m\cap(L^\times)^{p^n}
                       \subset U_m^{p^{n-m}}.                 \tag{META.AD2}
\]
The induced $p$-power topology is therefore the complete
logarithmic topology.  Thus $U_m$ injects into
$\varprojlim_n L^\times/(L^\times)^{p^n}$.
The inverse-limit Kummer sequence identifies this
limit with $H^1(G_L,\mathbb Z_p(1))$; the possible
$\varprojlim^1$ term vanishes since every group
$\mu_{p^n}(L)$ is finite and these inverse systems
are Mittag--Leffler.
Continuous cocycles in $\mathbb Q_p(1)$ have compact,
hence bounded, image, so scaling makes them integral.
The same observation for coboundaries proves
$H^1(G_L,\mathbb Z_p(1))\otimes\mathbb Q_p
=H^1(G_L,\mathbb Q_p(1))$.
Tensoring the injection and using logarithm gives
(META.AD1).  Its displayed formula proves independence
of $r$, compatibility with finite extensions, and
Galois equivariance.
In particular
\[
 v(u-1)>\frac1{p-1},\quad \kappa_L(u)=0
                         \quad\Longrightarrow\quad u=1.       \tag{META.AD3}
\]
No assertion about Kummer classes of arbitrary residue
units has entered this argument.

For completeness the inertia comparison also does not
use a residue-field restriction.
Embed $\overline{\mathbb Q}_p$ into $\overline M$ over
the canonical $\mathbb Q_p\subset M$, and put
$T=M^{\mathrm{nr}}\cap\overline{\mathbb Q}_p$.
Unramified constant extensions are finite etale after
base change to the valuation ring of $M$; hence
$\mathbb Q_p^{\mathrm{nr}}\subset T$.
If $A/\mathbb Q_p$ is finite inside $T$ and
$e=v_M(p)$ in integer normalization, its ramification
index divides $e$, since $M^{\mathrm{nr}}$ has the same
value group as $M$.
An Eisenstein uniformizer polynomial after the maximal
unramified subextension proves
\[
 [A\mathbb Q_p^{\mathrm{nr}}:\mathbb Q_p^{\mathrm{nr}}]
                                      =e(A/\mathbb Q_p)\le e.
\]
Thus $T/\mathbb Q_p^{\mathrm{nr}}$ has bounded finite
subextension degrees and is finite; write
$T=L_0\mathbb Q_p^{\mathrm{nr}}$ with $L_0/\mathbb Q_p$ finite.
The image of $I_M$ is the closed group with fixed field
$T$, namely $I_{L_0}$.
Its cyclotomic image is infinite, so
$\mathbb Q_p(1)^{I_M}=0$ and inflation--restriction
gives the claimed injectivity.
Iwasawa's Theorem~2(ii) says that the wild inertia of
$L_0$ is free pro-$p$ of infinite rank
\citep{AN_Iwasawa1955}.
Mapping finitely many generators to adjacent elementary
matrices in $\operatorname{UT}_{2^s+1}(\mathbb F_p)$,
and using
$[1+E_{ab},1+E_{bc}]=1+E_{ac}$ for $a<b<c$,
shows that every derived subgroup of that wild inertia
is nontrivial.
From the first commutator onward the derived groups
of $I_{L_0}$ are pro-$p$.  A nontrivial pro-$p$ group
has a nontrivial abelian quotient, as is seen from a
nontrivial finite quotient and then a quotient of order
$p$ of its abelianization.
The tame quotient treats the first derived step.
Surjectivity $I_M\to I_{L_0}$ carries closed derived
subgroups onto the corresponding closed derived
subgroups, proving strictness over $M$.

We now prove the evaluation detector.
The case $n=0$ is the asserted zero evaluation itself;
assume $n>0$.
After a finite scalar extension split the prime
boundary equations of $V$, say $h_1,\ldots,h_s$.
The geometric coordinate ring is a localization of
a polynomial ring, so has zero Picard group and has
units modulo constants freely generated by these
distinct prime equations.
Kummer theory and inflation--restriction consequently
give
\[
              \eta=\beta+\sum_{j=1}^s\lambda_j\kappa(h_j),
                         \lambda_j\in\mathbb Q_p.             \tag{META.AD4}
\]
Choose one algebraic parameter $t_0$ in the given ball.
Choose a finite Galois extension $M'/M_0$ containing
$L$, all coefficients of the $h_j$, and generators of
the finite scalar fields just used.
By the completion comparison, its algebraic part
$M'\cap\Omega$ is finite over $M_0\cap\Omega$.
Choose a finite extension $E'/E_0$ which contains $E$,
these coefficients, and $M'\cap\Omega$, with
\[
                       G_{E'}=N\cap G_{M'}=:H.
\]
Concretely take the compositum of $E_0$ with that
algebraic part, after enlarging $M'$ to include the
finitely many specified generators.  This is a finite
extension of $E_0$, and also of $E$.
The zero evaluation at $t_0$ says over $E'$ that
\[
                  \beta=-\sum_j\lambda_j\kappa(h_j(t_0)).
\]
For every actual parameter $t$ in a smaller ball,
all ratios $h_j(t)/h_j(t_0)$ lie in a deep unit group
of $M'$, and the local class
\[
 \theta_t=\sum_j\lambda_j
                   \kappa_{M'}(h_j(t)/h_j(t_0))
          =\iota_{M'}\left(
                \sum_j\lambda_j\log\frac{h_j(t)}{h_j(t_0)}
                         \right)                              \tag{META.AD5}
\]
restricts to zero on $H$.
Choose $g\in N$ whose restriction to $L$ is nontrivial.
It is essential that this is the original $N$, not the
smaller group $H$ which fixes $L$.
Since $M'/M_0$ is Galois and $N$ is normal in $G_{M_0}$,
$g$ normalizes $G_{M'}$ and
$[g,G_{M'}]\subset N\cap G_{M'}=H$.
It acts trivially on $\mathbb Q_p(1)$ as well.
Inflation--restriction makes $\theta_t$ a class from
$H^1(G_{M'}/H,\mathbb Q_p(1))$; $g$ acts as the identity
on this group, and hence $g\theta_t=\theta_t$.
Injectivity and equivariance of (META.AD1) give
\[
 \Phi(t)=g\Phi(t),\qquad
            \Phi(t)=\sum_j\lambda_j\log\frac{h_j(t)}{h_j(t_0)}.
                                                               \tag{META.AD6}
\]
No coefficient of $h_j$ or coordinate of $t_0$ is being
assumed fixed by $g$.

Here is the analytic consequence with those coefficient
actions retained.  Choose an $M_0$-basis $e_1,\ldots,e_q$
of $L$ lying in $L\cap\Omega$, and write, in each of the
$n$ coordinates,
$t=t_0+\sum_{\nu=1}^q u_\nu e_\nu$.
The algebraic scalar parameters $u_\nu\in M_0\cap\Omega$
are dense in a scalar $M_0$-ball.
After extending coefficients to $M'$, the invertible
embedding matrix $(\sigma(e_\nu))_{\sigma,\nu}$ gives
independent embedding blocks $z_\sigma$.
The left side of (META.AD6) is an analytic function
only of the identity block; the right side is the
analytic function with conjugated coefficients, centred
at $g t_0$, only of the distinct block $\sigma=g|_L$.
Equality on the dense algebraic scalar set gives an
identity of convergent germs.  To justify this without
local compactness, a nonzero one-variable convergent
power series cannot vanish on a scalar open ball, by
its first nonzero term near an accumulation point;
induction on the variables gives the same assertion
over a finite coefficient extension.
Differentiating in the identity block therefore gives
\[
                         \sum_j\lambda_j\,dh_j/h_j=0.
                                                               \tag{META.AD7}
\]
This identity of germs is an identity of rational
differential forms.  The residue at the distinct prime
divisor $h_j=0$ is $\lambda_j$, so every $\lambda_j$ is
zero.  The normalized arithmetic term is then zero as
well, proving $\eta|_{E'}=0$.
Restriction to a finite scalar extension is injective
with rational coefficients: its composite with
corestriction is multiplication by the nonzero degree.
Consequently $\eta=0$ as asserted.
\end{proof}

\begin{claim}{B.META.23}{Open}{Draft: reconstruction at derived and nilpotent inertia fields for arbitrary discrete mixed-characteristic bases}
Keep $D,M,\Omega$ of Claim~B.META.22 and put $I=I_M$.
Define
\[
 F_r=\Omega^{I^{(r)}}\quad(r\ge0),\qquad
 K_{\rm nil}=\Omega^{\bigcap_{c\ge1}\gamma_{c+1}(I)}.
                                                               \tag{META.AD8}
\]
Every finite extension of each $F_r$ and of $K_{\rm nil}$
satisfies full closed-point-preserving relative etale
reconstruction for finite-type schemes localized at
universal homeomorphisms, with the structural
isomorphism retained.  The property descends to any
subfield over which one of these fields is algebraic.
No residue-field finiteness hypothesis is imposed.
\end{claim}
\begin{proof}
This proof is being written after the checkpoint05
freeze.  The arithmetic argument and the local-to-global
norm argument below are separate from the frozen
Claims~B.META.20--B.META.21.

\medskip\noindent\textbf{Stable local comparisons, not a global kernel classification.}
For $r\ge2$ let $K=F_r$, $B=I^{(r-1)}$ and
$A=\Omega^B$.
If $E/K$ is finite and $\kappa_E(a)=0$, choose finite
Galois $D_1/D$ containing $a$ and generators of $E/K$.
Put $M_1$ for its completion, $J=I_{M_1}$,
$H=I^{(r)}\cap J$, $\Gamma=J/H$ and
$C=(B\cap J)/H$.
The local Kummer class inflates from $H^1(\Gamma,V)$,
where $V=\mathbb Q_p(1)$.
For $g\in B$, normality and Galoisness of $D_1/D$
make its actions on $C$, on $\Gamma/C$, and on $V$
the identity.  The inflation--restriction filtration
\[
 0\longrightarrow H^1(\Gamma/C,V)
 \longrightarrow H^1(\Gamma,V)
 \longrightarrow H^1(C,V)^{\Gamma/C}
\]
therefore gives $(g-1)^2=0$.
The particular Kummer vector has finite $g$-orbit
because $a\in D_1$.  If $g^m$ fixes it, then
$0=(g^m-1)\theta=m(g-1)\theta$ and $g\theta=\theta$.
Inertia restriction injectivity in Claim~B.META.22
now gives
\[
 \kappa_E(a)=0\quad\Longrightarrow\quad
             \kappa_{M_1}(g(a)/a)=0\qquad(g\in B).              \tag{META.AD9}
\]
We do not turn this into a root-of-unity assertion.
Instead, if $v(a-1)>1/(p-1)$, then $g(a)/a$ is in
the same deep neighbourhood of one and (META.AD3)
implies $g(a)=a$.

For $K=K_{\rm nil}$ put $N=\bigcap_c\gamma_{c+1}(I)$.
The arbitrary-DVR inertia theorem, Stacks Tag~0BUA,
gives a pro-$p$ wild group and the full prime-to-$p$
cyclic tame quotient; its proof makes no residue
perfectness assumption \citep{AN_Stacks0BUA}.
In finite nilpotent quotients all Sylow groups commute.
All tame quotients survive, and therefore
\[
              I/N=Q_{(p)}\times
                         \prod_{\ell\ne p}\mathbb Z_\ell(1),
                                                               \tag{META.AD10}
\]
where $Q_{(p)}$ is pro-$p$.
Choose an odd prime $\ell\nmid e p(p-1)$, where
$e=v_M(p)$ in integer normalization.
The element $b_0=p^{1/\ell}$ belongs to $K$ and
generates a degree-$\ell$ tame extension of the
strict henselization.  A lift $\tau\in I$ of a
central element in the pro-$\ell$ factor can be
chosen with $\tau b_0\ne b_0$.
Its $p$-cyclotomic value is one, since $\ell\nmid p(p-1)$;
inertia fixes all roots of unity of order prime to $p$.

If $\kappa_E(a)=0$ for finite $E/K$, the integral
Kummer class is already zero: $G_E$ fixes $\mathbb Z_p(1)$,
so its first cohomology is a torsion-free group of
continuous homomorphisms.  Thus compatible $p$-power
roots of $a$ lie in $E$.
Choose finite Galois $D_1/D$ containing the data and
replace $E$ by $KD_1$.
Let $F$ be the compositum of $D_1$, the relative
strict henselization, and $\mu_{p^\infty}$, and let
$R/F$ be generated by all $p$-power roots of the
finitely many inertia conjugates of $a$.
Then $R\subset KD_1$ and
$D_R=\operatorname{Gal}(R/F)$ is a torsion-free
closed subgroup of $\mathbb Z_p(1)^s$.
Both fields are Galois over the strict henselization.
The image of $N$ on $R$ is finite since $KD_1/K$ is
finite.  Centrality of $\tau$ modulo $N$ puts
$[\tau,D_R]$ in that finite image and in $D_R$;
thus it is trivial.
For $J=I_{M_1}$ and $J_0=J\cap\ker\chi_p$, the group
$J_0$ maps onto $D_R$, so this equality gives
$\kappa_{J_0}(\tau a)=\kappa_{J_0}(a)$.
Restriction from $H^1(J,V)$ to $H^1(J_0,V)$ is
injective: $J/J_0$ is abelian with nontrivial
cyclotomic action, and comparing a cocycle at $xy$
and $yx$ shows that it is a coboundary.
Inertia restriction injectivity finally yields
\[
 \kappa_E(a)=0\quad\Longrightarrow\quad
                         \kappa_{M_1}(\tau(a)/a)=0.            \tag{META.AD11}
\]
Again (META.AD3) makes $\tau(a)=a$ whenever $a$ is
in the deep neighbourhood of one.

\medskip\noindent\textbf{Nonzero actual classes and a local two-grid test.}
In the derived case strictness gives $b_0\in K\setminus A$;
choose $g\in B$ moving it.
In the nilpotent case use $b_0$ and $g=\tau$ above.
Replace $b_0$ by $b=p^m b_0$ for $m$ sufficiently large
that $v(b)>1/(p-1)$.
If $\kappa_K(1-b)$ were zero, (META.AD9) or (META.AD11)
and (META.AD3) would give
$(1-gb)/(1-b)=1$, contrary to $gb\ne b$.
Thus $\kappa_K(1-b)\ne0$.
Rational restriction and corestriction show that this
class remains nonzero over every finite extension of $K$.
All these fields contain $\mu_{p^\infty}$.

For an integer $d\ge1$ set
\[
 \mathcal A_0=\{0,1,\ldots,4d\},\qquad
 \mathcal A_1=b+\mathcal A_0,\qquad
 \mathcal A=\mathcal A_0\cup\mathcal A_1.                       \tag{META.AD12}
\]
Let $R=P/Q$ be a quotient of monic degree-$d$ polynomials
over any finite extension of $K$.
Suppose all $R(a)$, $a\in\mathcal A$, are defined,
have zero Kummer class there, and satisfy
$v(R(a)-1)>1/(p-1)$.
For the chosen $g$ moving $b$, (META.AD9) or (META.AD11)
and (META.AD3) give $g(R(a))=R(a)$.
At first-grid centres $g$ fixes $a$, so $gR=R$ by
polynomial interpolation: the cross-multiplied numerator
has degree at most $2d-1$.
At translated centres this gives
$R(gb+j)=R(b+j)$ for $0\le j\le4d$.
Interpolation again gives a nonzero additive period
$gb-b$.  A rational function in characteristic zero
with such a period is constant, since its finite pole
set is invariant under translation and its polynomial
part has a nonzero finite difference unless constant.
Monicity gives $R(\infty)=1$, so $R=1$.
This is a local test at the identity; it is not being
asserted for arbitrary values in the full Kummer kernel.

% The local norm construction, principal-open detection,
% all-centres globalization, graph proof and strict example
% are being appended.  This proof is not yet complete.
\end{proof}
